\documentclass{scrartcl}
\usepackage{graphicx}  % required for inserting images

% for cyrillic and Bulgarian language
\usepackage[utf8]{inputenc}
\usepackage[T2A]{fontenc}
\usepackage[bulgarian]{babel}

% math packages
\usepackage{amsmath}
\usepackage{float}
\usepackage{amsfonts}
\usepackage{amssymb}
\usepackage{amsthm}
\usepackage{cancel}

\usepackage{ragged2e}  % adds FlushLeft

\usepackage{listings}
\usepackage{xcolor}

\definecolor{codegreen}{rgb}{0,0.6,0}
\definecolor{codegray}{rgb}{0.5,0.5,0.5}
\definecolor{codepurple}{rgb}{0.58,0,0.82}
\definecolor{backcolour}{rgb}{0.95,0.95,0.92}

\lstdefinestyle{mystyle}{
    %caption=Example C++,
    label={lst:listing-cpp},
    language=C++,
    backgroundcolor=\color{backcolour},   
    commentstyle=\color{codegreen},
    keywordstyle=\color{magenta},
    numberstyle=\tiny\color{codegray},
    stringstyle=\color{codepurple},
    basicstyle=\ttfamily\footnotesize,
    breakatwhitespace=false,         
    breaklines=true,                 
    keepspaces=true,                 
    numbers=left,       
    numbersep=5pt,                  
    showspaces=false,                
    showstringspaces=false,
    showtabs=false,                  
    tabsize=2,
}
\lstset{style=mystyle}

\usepackage[colorlinks=true, linkcolor=blue, urlcolor=cyan]{hyperref}

\newtheorem{theorem}{Теорема}
\newtheorem{definition}{Дефиниция}

\title{Обектно ориентирано програмиране - Полиморфизъм. Множествено наследяване.}
\author{Ивайло Андреев + Claude Sonnet 4.6}
\date{17 май 2026}

\begin{document}
\maketitle

\tableofcontents

\newpage

\begin{FlushLeft}

% ============================================================
\section{Полиморфизъм -- общ преглед}
% ============================================================

\textbf{Полиморфизмът} (от гр. „много форми") е принципът, при който един и същ интерфейс се реализира по различен начин в зависимост от типа на обекта. В C++ съществуват два вида:

\begin{itemize}
    \item \textbf{Статичен (compile-time) полиморфизъм:} изборът на функция се прави по време на компилация. Реализира се чрез претоварване на функции/оператори и \textbf{шаблони (templates)}.
    \item \textbf{Динамичен (run-time) / Подтипов полиморфизъм:} изборът се прави по време на изпълнение. Реализира се чрез \textbf{наследяване и виртуални функции}.
\end{itemize}

% ============================================================
\section{Виртуални функции и подтипов полиморфизъм. Динамично свързване.}
% ============================================================

\subsection{Статично срещу динамично свързване}

Без ключовата дума \texttt{virtual} компилаторът определя коя функция да извика единствено на базата на \textbf{статичния тип} на указателя/референцията (ранно свързване). С \texttt{virtual} -- на базата на \textbf{реалния тип} на обекта по време на изпълнение (късно свързване).

\begin{lstlisting}
class Base {
public:
    void f() const { std::cout << "Base::f()\n"; }       // non-virtual
    virtual void g() const { std::cout << "Base::g()\n"; } // virtual
};

class Derived : public Base {
public:
    void f() const { std::cout << "Derived::f()\n"; }          // hides Base::f
    void g() const override { std::cout << "Derived::g()\n"; } // overrides Base::g
};

int main() {
    Derived d;
    Base* p = &d;

    p->f(); // Base::f()    -- static binding (non-virtual)
    p->g(); // Derived::g() -- dynamic binding (virtual)
}
\end{lstlisting}

\subsection{Виртуални функции -- правила и спецификатори}

\begin{itemize}
    \item \textbf{\texttt{virtual}:} Декларира функцията за виртуална в базовия клас. Всички класове-наследници наследяват тази виртуалност автоматично.
    \item \textbf{\texttt{override}:} Изрично указва, че функцията предефинира виртуална функция от базовия клас. Компилаторът проверява -- предпазва от typo-грешки в сигнатурата.
    \item \textbf{\texttt{final}:} Забранява по-нататъшно предефиниране на функцията (или наследяване от класа).
    \item \textbf{Виртуален деструктор:} Задължителен в базов клас, ако обектите ще се изтриват чрез указател към базовия клас. Без него \texttt{delete p} извиква само деструктора на базовия клас -- \textbf{изтичане на памет} или неопределено поведение.
\end{itemize}

\begin{lstlisting}
class Shape {
public:
    virtual void draw() const { std::cout << "Drawing shape\n"; }
    virtual ~Shape() = default;       // virtual destructor -- mandatory!
};

class Circle : public Shape {
public:
    void draw() const override final { // override + forbid further overriding
        std::cout << "Drawing circle\n";
    }
};
\end{lstlisting}

\textbf{Ограничения:} Виртуалните функции не могат да бъдат \texttt{static} (нямат \texttt{this} указател). Конструкторите не могат да бъдат виртуални.

\subsection{Пример: динамично свързване при дълбока йерархия}

Следният пример илюстрира точно коя функция се извиква при различни комбинации от йерархия и override-ване:

\begin{lstlisting}
class Base {
public:
    virtual void f() const { std::cout << "Base::f()\n"; }
    virtual void g() const { std::cout << "Base::g()\n"; }
    void nonVirtual() const { std::cout << "Base::nonVirtual()\n"; }
};

class FirstDerived : public Base {
public:
    void f() const override { std::cout << "FirstDerived::f()\n"; }
    void g() const override { std::cout << "FirstDerived::g()\n"; }
    virtual void h() const  { std::cout << "FirstDerived::h()\n"; } // new virtual
};

class SecondDerived : public FirstDerived {
public:
    void f() const override { std::cout << "SecondDerived::f()\n"; }
    // g() NOT overridden -- inherited from FirstDerived
};

class ThirdDerived : public SecondDerived {
public:
    void h() const override { std::cout << "ThirdDerived::h()\n"; }
    // f() NOT overridden -- inherited from SecondDerived
};
\end{lstlisting}

\begin{lstlisting}
int main() {
    Base          baseObj;
    FirstDerived  firstObj;
    SecondDerived secondObj;
    ThirdDerived  thirdObj;

    Base* p = nullptr;

    p = &baseObj;
    p->f();          // Base::f()
    p->g();          // Base::g()
    p->nonVirtual(); // Base::nonVirtual()  (static -- always Base)

    p = &firstObj;
    p->f();          // FirstDerived::f()
    p->g();          // FirstDerived::g()

    p = &secondObj;
    p->f();          // SecondDerived::f()
    p->g();          // FirstDerived::g()  -- not overridden in SecondDerived!

    p = &thirdObj;
    p->f();          // SecondDerived::f() -- not overridden in ThirdDerived
    p->g();          // FirstDerived::g()  -- not overridden anywhere below First

    // h() is NOT in Base's interface -- cannot call through Base*
    // Need FirstDerived* to access h():
    FirstDerived* fp = &thirdObj;
    fp->h();         // ThirdDerived::h() -- dynamic dispatch
}
\end{lstlisting}

Ключовото наблюдение: виртуалната функция \textbf{продължава да се търси нагоре по йерархията} до намерения override. Ако \texttt{SecondDerived} не предефинира \texttt{g()}, но \texttt{FirstDerived} го е направил -- \texttt{FirstDerived::g()} се извиква. Невиртуалните методи (\texttt{nonVirtual()}) винаги се свързват статично с типа на указателя.

\subsection{Абстрактни класове и чисто виртуални функции}

\textbf{Чисто виртуална функция} се декларира с \texttt{= 0} и задължава всеки конкретен наследник да я предефинира. Клас с поне една чисто виртуална функция е \textbf{абстрактен} -- от него не могат да се създават обекти директно. Служи като интерфейс.

\begin{lstlisting}
class Animal {
public:
    virtual void speak() const = 0;  // pure virtual -- no body here
    virtual ~Animal() = default;
};

class Dog : public Animal {
public:
    void speak() const override { std::cout << "Woof!\n"; }
};

class Cat : public Animal {
public:
    void speak() const override { std::cout << "Meow!\n"; }
};

// Animal a;  // ERROR: cannot instantiate abstract class
Animal* a = new Dog();
a->speak();   // Woof! -- dynamic dispatch
delete a;
\end{lstlisting}

Чисто виртуална функция \textbf{може да има тяло}, дефинирано извън класа. Наследниците могат да го извикват явно чрез \texttt{Base::f()}. Класът остава абстрактен -- наследниците пак трябва да я предефинират.

\begin{lstlisting}
void Animal::speak() const {         // body of pure virtual
    std::cout << "...(silence)...\n";
}

class Parrot : public Animal {
public:
    void speak() const override {
        Animal::speak();             // call base body explicitly
        std::cout << "Polly wants a cracker!\n";
    }
};
\end{lstlisting}

% ============================================================
\section{Масиви от обекти и от указатели към обекти}
% ============================================================

При работа с наследяване се срещат четири основни варианта на деклариране:

\begin{lstlisting}
class Base    { public: virtual void f(); virtual ~Base() = default; };
class Derived : public Base { public: void f() override; };

Base     arr[10];   // array of 10 Base objects by value -- no polymorphism
Derived  arr[10];   // array of 10 Derived objects by value -- no polymorphism
Base*    arr;       // pointer to Base -- supports polymorphism (can point to Derived)
Derived* arr;       // pointer to Derived -- Derived objects only
\end{lstlisting}

Стойностните масиви (\texttt{Base arr[]} и \texttt{Derived arr[]}) \textbf{не поддържат полиморфизъм} -- при съхранение на производен обект в Base слот настъпва \textbf{object slicing} (виж 3.4). Само масивите от указатели (\texttt{Base* arr[]}) позволяват динамично свързване.

\subsection{Масив от обекти (на стека)}

\begin{lstlisting}
class Car { /* ... */ };

Car arr[5];   // calls default constructor 5 times
              // destructors called automatically when arr goes out of scope
\end{lstlisting}

\textbf{Характеристики:} Обектите са плътно в паметта. Всички елементи се конструират с един и същ конструктор (по подразбиране). Размерът е фиксиран при компилация.

\subsection{Динамичен масив от обекти (на heap-а)}

\begin{lstlisting}
Car* arr = new Car[5];  // calls default constructor 5 times
delete[] arr;           // calls destructor 5 times, then frees memory
                        // delete[] -- NOT delete!
\end{lstlisting}

Размерът може да се определи по време на изпълнение. Изисква \texttt{delete[]} -- \texttt{delete} (без \texttt{[]}) е \textbf{неопределено поведение}.

\subsection{Масив от указатели към обекти}

\begin{lstlisting}
Car* arr[5];               // array of 5 pointers (uninitialized)
arr[0] = new Car(42);      // each element independently constructed
arr[1] = new Car(7);
// ...
for (int i = 0; i < 5; i++) delete arr[i]; // each element independently destroyed
\end{lstlisting}

\textbf{Предимства:} всеки елемент може да е конструиран с различни аргументи; позиция може да бъде \texttt{nullptr}; позволява \textbf{полиморфизъм}.

\subsection{Проблемът с нарязването (Object Slicing) и полиморфните контейнери}

При съхранение на производен обект \textbf{по стойност} в базов тип се губят допълнителните полета и виртуалното поведение -- това се нарича \textbf{object slicing}:

\begin{lstlisting}
class Base {
public:
    virtual std::string name() const { return "Base"; }
    virtual ~Base() = default;
};

class Derived : public Base {
public:
    std::string name() const override { return "Derived"; }
    int extra = 42;  // extra data, lost on slicing
};

// Problem: storing by value causes slicing
Derived d;
Base b = d;      // SLICING: only the Base part is copied
b.name();        // "Base" -- NOT "Derived"!
\end{lstlisting}

Решението е \textbf{масив от указатели към базовия клас} -- полиморфен контейнер:

\begin{lstlisting}
Base* arr[2];
arr[0] = new Base();
arr[1] = new Derived();

for (int i = 0; i < 2; i++) {
    std::cout << arr[i]->name() << "\n"; // "Base", then "Derived" -- correct!
    delete arr[i];                        // correct because ~Base() is virtual
}
\end{lstlisting}

\textbf{Задължително:} Деструкторът на базовия клас трябва да е \texttt{virtual}, за да се извикат правилните деструктори на наследниците при \texttt{delete arr[i]}.

% ============================================================
\section{Параметричен полиморфизъм. Шаблони на функция и на клас.}
% ============================================================

Параметричният полиморфизъм се реализира чрез \textbf{шаблони (templates)} -- механизъм за писане на обобщен код, работещ с произволен тип. Типът е параметър, подаден при компилация. Компилаторът \textbf{инстанцира} (генерира конкретен код) за всеки използван тип.

\subsection{Шаблони на функция}

\begin{lstlisting}
template <typename T>
T maxOf(const T& a, const T& b) {
    return (a > b) ? a : b;
}

int main() {
    std::cout << maxOf<int>(3, 7);      // explicit: T = int
    std::cout << maxOf(3.14, 2.71);     // deduced: T = double (type deduction)
    std::cout << maxOf(std::string("a"), std::string("b")); // T = string
}
\end{lstlisting}

Компилаторът генерира отделна функция за всеки тип \texttt{T} -- \texttt{maxOf<int>}, \texttt{maxOf<double>} и т.н. Ако типът не поддържа \texttt{>}, се получава грешка при компилация.

Може да се използва \texttt{class} вместо \texttt{typename} -- двете са еквивалентни в този контекст.

\subsection{Шаблони на клас}

\begin{lstlisting}
template <typename T>
class Array {
public:
    Array(int n) : size_(n), data_(new T[n]) { }
    ~Array() { delete[] data_; }
    T& operator[](int i) { return data_[i]; }
    int size() const     { return size_; }
private:
    int  size_;
    T*   data_;
};

int main() {
    Array<int> nums(5);
    nums[0] = 10;
    std::cout << nums.size(); // 5

    Array<std::string> words(3);
    words[0] = "hello";
}
\end{lstlisting}

\subsection{Нетипови параметри на шаблон}

\texttt{Array<T>} от предишната секция използва динамична памет (heap). Параметрите на шаблона могат да бъдат и конкретни стойности (\textbf{нетипови параметри}), което позволява размерът да е известен при компилация -- тогава масивът е на стека, без heap заделяне. \texttt{Array<T, N>} по-долу е \textbf{различна дефиниция}, не предефиниране на същата.

Параметрите на шаблон могат да бъдат и конкретни стойности (не само типове):

\begin{lstlisting}
template <typename T, int N>
class Array {
    T data[N];   // fixed-size array, size known at compile time
public:
    int size() const { return N; }
    T& operator[](int i) { return data[i]; }
};

Array<double, 10> v;  // array of 10 doubles, no heap allocation
\end{lstlisting}

\subsection{Специализация на шаблон}

Можем да дефинираме специално поведение за конкретен тип:

\begin{lstlisting}
template <typename T>
T sum(T a, T b) { return a + b; }   // general

template <>
const char* sum(const char* a, const char* b) {  // specialization for char*
    // custom string concatenation logic
}
\end{lstlisting}

\textbf{Специализация от стандартната библиотека:} \texttt{std::vector<bool>} е специализация на \texttt{std::vector<T>} -- вместо масив от \texttt{bool} стойности, елементите се пакетират побитово (\textbf{bitset оптимизация по памет}: 1 бит на елемент вместо 1 байт). Заради тази оптимизация обаче \texttt{operator[]} не връща \texttt{bool\&}, а proxy обект, което нарушава стандартните очаквания за контейнер и е считано за проектна грешка.

\subsection{Ограничение: дефиниции в хедър файл}

Шаблоните не могат да се разделят на \texttt{.h} (декларация) и \texttt{.cpp} (дефиниция) по традиционния начин. При инстанциране компилаторът трябва да вижда цялото тяло на шаблона. Ако дефиницията е в отделен \texttt{.cpp}, linker-ът ще съобщи за липсващ символ. \textbf{Стандартна практика:} цялата имплементация в хедър файла (\texttt{.h}).

% ============================================================
\section{Множествено наследяване}
% ============================================================

C++ позволява производен клас да има \textbf{повече от един директен базов клас}:

\begin{lstlisting}
class Flyable {
public:
    virtual void fly() const { std::cout << "Flying\n"; }
};

class Swimmable {
public:
    virtual void swim() const { std::cout << "Swimming\n"; }
};

class Duck : public Flyable, public Swimmable {
public:
    void quack() const { std::cout << "Quack!\n"; }
};

int main() {
    Duck d;
    d.fly();   // OK
    d.swim();  // OK
    d.quack(); // OK
}
\end{lstlisting}

\textbf{Предимства:} комбиниране на независими интерфейси/функционалности; повторна употреба на код.

\subsection{Диамантен проблем (Diamond Problem)}

Проблемът възниква, когато един базов клас се наследява по \textbf{два независими пътя}, което води до \textbf{две отделни копия} на базовия клас в производния.

\begin{lstlisting}
class Base { public: int x; };
class A : public Base { };    // A::Base::x
class B : public Base { };    // B::Base::x -- separate copy!
class C : public A, public B { };

int main() {
    C obj;
    // obj.x = 5;       // ERROR: ambiguous -- A::x or B::x ?
    obj.A::x = 5;       // OK: explicit disambiguation
    obj.B::x = 10;      // OK
}
\end{lstlisting}

Конструкторът на \texttt{Base} се извиква \textbf{два пъти} -- веднъж чрез \texttt{A} и веднъж чрез \texttt{B}.

\subsection{Виртуално наследяване -- решение на диамантения проблем}

Ключовата дума \texttt{virtual} пред базовия клас гарантира, че в крайния наследник ще съществува \textbf{само едно споделено копие} на общия базов клас.

\begin{lstlisting}
class Base { public: int x; };
class A : virtual public Base { };   // virtual inheritance
class B : virtual public Base { };   // virtual inheritance
class C : public A, public B { };    // one shared Base::x

int main() {
    C obj;
    obj.x = 42;  // OK: no ambiguity, only one Base
}
\end{lstlisting}

\subsection{Конструктори при виртуално наследяване}

При виртуално наследяване \textbf{само най-производният клас} (в случая \texttt{C}) инициализира виртуалния базов клас. Инициализациите на \texttt{Base} в \texttt{A} и \texttt{B} се \textbf{игнорират}:

\begin{lstlisting}
class Base {
public:
    Base(int v) { std::cout << "Base(" << v << ")\n"; }
};

class A : virtual public Base {
public:
    A() : Base(0) { } // Base(0) is IGNORED when constructing through C
};

class B : virtual public Base {
public:
    B() : Base(1) { } // Base(1) is IGNORED when constructing through C
};

class C : public A, public B {
public:
    C() : Base(42), A(), B() { } // Base(42) is the one that runs
};
// Output: Base(42)  A()  B()  C()
\end{lstlisting}

\textbf{Ред на конструиране при виртуално наследяване:}
\begin{enumerate}
    \item Конструктори на виртуалните базови класове (извикани от най-производния клас).
    \item Конструктори на невиртуалните базови класове (в реда на декларирането им).
    \item Конструктори на член-данните (в реда на декларирането им).
    \item Тялото на конструктора на самия клас.
\end{enumerate}

\end{FlushLeft}

\end{document}